\section{Glossary}\label{supp:glossary}
Plain-language meanings of the terms the main paper uses most. The definitions in the main paper govern where the two differ.

{\small
\begin{description}[leftmargin=0pt,itemsep=2pt]
\item[Accountability floor ($a_{\min}$).] The share of revenue that a fully automated firm still spends on being accountable: officer certifications, the integrated audit, internal control over financial reporting, insurance and legal costs. Part of it is a requirement on how controls are designed, not a headcount (\autoref{prop:amin}).
\item[Anchor.] The firm-years at or above the 95th percentile of the 2023--25 cash score, without acquisition-affected years, from which the central path starts (Section~\ref{sec:empirics-anchor}; Table~\ref{tab:anchor-firms}).
\item[Capital intensity ($k$).] Capital the firm must add per dollar of new revenue. With $k>0$ growth uses cash; with $k<0$, as when customers pay in advance, growth brings cash in.
\item[Compute share ($c_c$).] The cost of the computing that replaces labor, as a share of revenue, at full automation. A scenario value, and the largest uncertainty in the ceiling.
\item[Decision weight.] The weight a value-maximizing firm puts on a point of growth relative to a point of margin: under the verification hypothesis (H) of \autoref{prop:dynamic}, the value of an extra dollar of revenue. It depends on the owner's discount rate, and the data do not identify it.
\item[Demand pull.] How much growth a dollar of growth spend buys at a small market share; better products and channels have more of it.
\item[Durable (saturated) score.] The score a firm keeps once it has no share left to gain: its margin intercept plus market growth net of the capital that growth needs.
\item[Embedded and realized shares.] Of an automation program, the share whose processes have been changed and stay changed (embedded), and the share of the program's targeted benefits the firm actually captures (realized), which is smaller.
\item[Estimands E1 and E2.] E1 uses scores as reported; E2 removes growth bought by acquisition. The imputation bounds add back, with imputed organic growth, the acquisition years E2 cannot adjust.
\item[Familywise correction.] An adjustment of $p$-values for testing many statistics at once, so that the chance of any false finding in the family stays at the stated level.
\item[Frontier.] The best margin a firm can earn at each growth rate. The frontier estimated across firms (\autoref{def:industry}) is drawn through the points firms chose, not through any one firm's possibilities.
\item[Full-automation ceiling ($R_{\max}$).] The highest durable cash score a firm can reach at a stated full-automation cost structure and fixed prices (\autoref{cor:ceiling}). A bound at that cost structure, not a forecast or a target.
\item[Full-cost floors.] The recurring costs of running AI that full automation adds: maintaining, securing and supervising it, and the losses its errors cause (\autoref{def:fullcost}).
\item[Growth surplus ($\Pi$).] The part of a score above its durable level, earned while the firm still has share to gain. It is transient.
\item[Margin intercept ($m_0$).] The margin a firm would earn if it spent nothing to grow faster than its market.
\item[Market growth ($G$).] How fast the firm's whole market grows; a firm can match it without growth spend.
\item[Pass-through.] The share of a fall in unit costs that competition passes to customers as lower prices.
\item[Price line.] A Rule-of-40 or Rule-of-$X$ line: its slope is the number of margin points the rule trades for one point of growth. Firms on the same line have the same score, not the same value.
\item[$\Psi_\chi$.] The paper's primary frontier statistic: for each growth level, the $\chi$-quantile of the margin among firm-years that grew at least that fast.
\item[Remote threshold ($\bar\pi$).] In the control model, the largest probability that a design may leave that, over the year's opportunities for a material misstatement, at least one passes every control layer. A scenario value, not a tolerance.
\item[Screening weight.] The weight on growth that investors use to compare firms, typically two to three; it reflects how long growth lasts.
\item[Share headroom and the score budget.] The room a firm has to gain share in its market. Within one market and at given retention, every point of score above the durable level uses up part of a finite budget of headroom (\autoref{prop:budget}).
\item[Spend elasticity ($\alpha$).] How fast the returns to growth spend diminish; the data do not identify it, so the paper uses a range.
\item[Top-tail firm.] A firm at the 95th percentile of the 2023--25 cash score, the starting point of the central path.
\item[$T(g_0)$.] The 95th-percentile margin among firm-years growing at least $g_0$, plus $g_0$: a point on the frontier at a fixed growth rate.
\end{description}
}
